\chapter{Apendice de Codigo Completo das Solucoes}
\label{apendice:codigo_completo}

Este apendice consolida o codigo completo dos artefatos discutidos ao longo dos capitulos.
O objetivo e manter a rastreabilidade entre explicacao e implementacao executavel no projeto.
Para estabilidade de compilacao em \LaTeX{} com \texttt{listings}, os blocos abaixo usam
copias ASCII geradas em \texttt{livro/generated/apendice\_ascii/}; os arquivos originais
permanecem no projeto com seus caminhos de origem.

\section{Como validar a autocontencao}

Antes de executar os exemplos, rode:

\begin{lstlisting}[language=bash,caption={Verificacao automatica de referencias de codigo no livro.}]
pwsh -File livro/validar_autocontencao.ps1
\end{lstlisting}

Resultado esperado: \texttt{Referencias faltantes: 0}.

% Ajuste local para melhorar quebra de linhas em arquivos extensos do apendice
\lstset{basicstyle=\ttfamily\footnotesize}

\section{Capitulo 2 e Capitulo 10}

\subsection{langchain.ipynb (exportado para Python)}
\lstinputlisting[language=Python,caption={Notebook base do laboratorio (codigo completo dos blocos, versao ASCII para compilacao).}]{generated/apendice_ascii/livro__generated__langchain_notebook_completo.py}

\subsection{\texttt{demo\_api.ipynb} (exportado para Python)}
\lstinputlisting[language=Python,caption={Notebook de validacao de API ponta a ponta (codigo completo dos blocos, versao ASCII para compilacao).}]{generated/apendice_ascii/livro__generated__demo_api_notebook_completo.py}

\section{Capitulo 4 e Capitulo 5}

\subsection{\texttt{mcp\_mock\_server.py}}
\lstinputlisting[language=Python,caption={Servidor MCP mock para testes locais (versao ASCII para compilacao).}]{generated/apendice_ascii/mcp_mock_server.py}

\subsection{\texttt{mcp\_postgres\_server.py}}
\lstinputlisting[language=Python,caption={Servidor MCP para PostgreSQL (versao ASCII para compilacao).}]{generated/apendice_ascii/mcp_postgres_server.py}

\subsection{\texttt{validate\_sprint\_a.py}}
\lstinputlisting[language=Python,caption={Script de validacao do Sprint A (versao ASCII para compilacao).}]{generated/apendice_ascii/validate_sprint_a.py}

\section{Capitulo 11}

\subsection{\texttt{otel\_config.py}}
\lstinputlisting[language=Python,caption={Configuracao de observabilidade (OTel/Langfuse, versao ASCII para compilacao).}]{generated/apendice_ascii/infraestrutura__modules__lambda__src__otel_config.py}

\subsection{status.py}
\lstinputlisting[language=Python,caption={Endpoint de consulta de status (versao ASCII para compilacao).}]{generated/apendice_ascii/infraestrutura__modules__lambda__src__status.py}

\subsection{worker.py}
\lstinputlisting[language=Python,caption={Worker principal do pipeline assincrono (versao ASCII para compilacao).}]{generated/apendice_ascii/infraestrutura__modules__lambda__src__worker.py}

\section{Capitulo 12, 13 e 14}

\subsection{\texttt{avaliar\_metricas.py}}
\lstinputlisting[language=Python,caption={Metricas automaticas de avaliacao (versao ASCII para compilacao).}]{generated/apendice_ascii/avaliar_metricas.py}

\subsection{\texttt{llm\_judge.py}}
\lstinputlisting[language=Python,caption={Avaliacao qualitativa com LLM as Judge (versao ASCII para compilacao).}]{generated/apendice_ascii/llm_judge.py}

\subsection{\texttt{golden\_dataset.py}}
\lstinputlisting[language=Python,caption={Criacao e regressao com golden dataset (versao ASCII para compilacao).}]{generated/apendice_ascii/golden_dataset.py}

\subsection{\texttt{golden\_baseline.json}}
\lstinputlisting[caption={Baseline de referencia para comparacao (versao ASCII para compilacao).}]{generated/apendice_ascii/golden_baseline.json}

\section{Capitulo 18}

\subsection{app.js}
\lstinputlisting[caption={Frontend de operacao e ingestao (versao ASCII para compilacao).}]{generated/apendice_ascii/frontend__cloudfront_console__app.js}

\subsection{config.js}
\lstinputlisting[caption={Configuracao do frontend CloudFront (versao ASCII para compilacao).}]{generated/apendice_ascii/frontend__cloudfront_console__config.js}
